% GENERATED FILE. Edit canonical lessons, then run npm run generate:books.
% canonical-source-hash: fnv1a64:5eb64def9b8b7a60
% canonical-lessons: PA-C261-addha, PA-C261-paletfaram, PA-C261-thahirna, PA-C261-challna, PA-C261-ravana

\chapter{The Train Timetable}
\label{ch:pa-a2-the-train-timetable}
\hlchaptermodality{261}

\begin{chapteropening}
\textbf{By the end of this chapter:} \emph{I can understand when a train leaves and where it stops.}

Complete the last lesson of chapter 261: I can understand when a train leaves and where it stops.
\end{chapteropening}

\section[\texorpdfstring{\pa{ਅੱਧਾ} (addhā)}{ਅੱਧਾ (addhā)}]{\pa{ਅੱਧਾ} (addhā) — half}

\label{lesson:PA-C261-addha}



\begin{quote}
Before the new one: say the Punjabi for one and a half, then the Punjabi for two and a half.
\end{quote}

\subsection*{You'll want to know: \pa{ਅੱਧਾ}}
\textbf{\pa{ਅੱਧਾ}} — \emph{addhā} — \textquotedblleft{}half\textquotedblright{}.

\begin{quote}
\pa{ਮੈਂ} \pa{ਅੱਧਾ} \pa{ਘੰਟਾ} \pa{ਸਟੇਸ਼ਨ} \pa{ਉੱਤੇ} \pa{ਉਡੀਕਦਾ} \pa{ਹਾਂ।}
\end{quote}

\begin{quote}
\emph{maiṁ addhā ghaṇṭā saṭeshan utte uḍīkdā hāṁ.}
\end{quote}

\begin{quote}
I wait half an hour at the station.
\end{quote}

\begin{cousinweb}
\textbf{\pa{ਅੱਧਾ}} (\emph{addhā}) is inherited from Sanskrit \emph{ardha}, \textquotedblleft{}half\textquotedblright{}; CDIAL lists Punjabi \emph{addhā}, \textquotedblleft{}half\textquotedblright{}, under it (CDIAL 644).
\end{cousinweb}

\begin{grammarlens}[title={What you've built}]
One more word for reading notices and messages.
\end{grammarlens}

\subsection*{Your turn}
\begin{itemize}
\raggedright
  \item \emph{Say it:} \emph{addhā}
  \item \emph{Say it:} \emph{addhā}, once more
  \item \emph{Say it:} \emph{maiṁ addhā ghaṇṭā saṭeshan utte uḍīkdā hāṁ.} — I wait half an hour at the station.
  \item \emph{Recall:} say the Punjabi for one and a half, then the Punjabi for two and a half, then say \emph{addhā} again
\end{itemize}

\begin{tcolorbox}[breakable,colback=teal!4,colframe=teal!35!black,title={Before you move on}]
What does \pa{ਅੱਧਾ} mean? (\textquotedblleft{}Half\textquotedblright{}.) Say it once more.
\end{tcolorbox}
\section[\texorpdfstring{\pa{ਪਲੇਟਫ਼ਾਰਮ} (paleṭfāram)}{ਪਲੇਟਫ਼ਾਰਮ (paleṭfāram)}]{\pa{ਪਲੇਟਫ਼ਾਰਮ} (paleṭfāram) — a railway platform}

\label{lesson:PA-C261-paletfaram}



\begin{quote}
Before the new one: say the Punjabi for two and a half, then the Punjabi for half.
\end{quote}

\subsection*{You'll want to know: \pa{ਪਲੇਟਫ਼ਾਰਮ}}
\textbf{\pa{ਪਲੇਟਫ਼ਾਰਮ}} — \emph{paleṭfāram} — \textquotedblleft{}a railway platform\textquotedblright{}.

\begin{quote}
\pa{ਰੇਲ} \pa{ਦੂਜੇ} \pa{ਪਲੇਟਫ਼ਾਰਮ} \pa{ਉੱਤੇ} \pa{ਆਵੇਗੀ।}
\end{quote}

\begin{quote}
\emph{rel dūje paleṭfāram utte āvegī.}
\end{quote}

\begin{quote}
The train will come in at the second platform.
\end{quote}

A masculine noun.

\begin{cousinweb}
\textbf{\pa{ਪਲੇਟਫ਼ਾਰਮ}} (\emph{paleṭfāram}) is borrowed from English \emph{platform}.
\end{cousinweb}

\begin{grammarlens}[title={What you've built}]
One more word for reading notices and messages.
\end{grammarlens}

\subsection*{Your turn}
\begin{itemize}
\raggedright
  \item \emph{Say it:} \emph{paleṭfāram}
  \item \emph{Say it:} \emph{paleṭfāram}, once more
  \item \emph{Say it:} \emph{rel dūje paleṭfāram utte āvegī.} — The train will come in at the second platform.
  \item \emph{Recall:} say the Punjabi for two and a half, then the Punjabi for half, then say \emph{paleṭfāram} again
  \item \emph{Write it:} copy \textbf{\pa{ਪਲੇਟਫ਼ਾਰਮ}} once, checking the whole word
\end{itemize}

\begin{tcolorbox}[breakable,colback=teal!4,colframe=teal!35!black,title={Before you move on}]
What does \pa{ਪਲੇਟਫ਼ਾਰਮ} mean? (\textquotedblleft{}A railway platform\textquotedblright{}.) Say it once more.
\end{tcolorbox}
\section[\texorpdfstring{\pa{ਠਹਿਰਨਾ} (ṭhahirnā)}{ਠਹਿਰਨਾ (ṭhahirnā)}]{\pa{ਠਹਿਰਨਾ} (ṭhahirnā) — to halt; to stop for a while}

\label{lesson:PA-C261-thahirna}



\begin{quote}
Before the new one: say the Punjabi for half, then the Punjabi for a railway platform.
\end{quote}

\subsection*{You'll want to know: \pa{ਠਹਿਰਨਾ}}
\textbf{\pa{ਠਹਿਰਨਾ}} — \emph{ṭhahirnā} — \textquotedblleft{}to halt; to stop for a while\textquotedblright{}.

\begin{quote}
\pa{ਰੇਲ} \pa{ਹਰ} \pa{ਸਟੇਸ਼ਨ} \pa{ਉੱਤੇ} \pa{ਦਸ} \pa{ਮਿੰਟ} \pa{ਠਹਿਰਦੀ} \pa{ਹੈ।}
\end{quote}

\begin{quote}
\emph{rel har saṭeshan utte das minṭ ṭhahirdī hai.}
\end{quote}

\begin{quote}
The train halts for ten minutes at every station.
\end{quote}

Past (it …): \pa{ਉਹ} \pa{ਰੇਲ} \pa{ਸਟੇਸ਼ਨ} \pa{ਉੱਤੇ} \pa{ਠਹਿਰੀ}, \emph{oh rel saṭeshan utte ṭhahirī}.

\begin{cousinweb}
\textbf{\pa{ਠਹਿਰਨਾ}} (\emph{ṭhahirnā}) is inherited from a reconstructed \emph{*stabhira}, \textquotedblleft{}fixed, firm\textquotedblright{}; CDIAL lists Punjabi \emph{ṭhahirnā}, \textquotedblleft{}to stop, be settled\textquotedblright{}, under it (CDIAL 13680).
\end{cousinweb}

\begin{grammarlens}[title={What you've built}]
One more word for reading notices and messages.
\end{grammarlens}

\subsection*{Your turn}
\begin{itemize}
\raggedright
  \item \emph{Say it:} \emph{ṭhahirnā}
  \item \emph{Say it:} \emph{ṭhahirnā}, once more
  \item \emph{Say it:} \emph{rel har saṭeshan utte das minṭ ṭhahirdī hai.} — The train halts for ten minutes at every station.
  \item \emph{Recall:} say the Punjabi for half, then the Punjabi for a railway platform, then say \emph{ṭhahirnā} again
  \item \emph{Write it:} copy \textbf{\pa{ਠਹਿਰਨਾ}} once, checking the whole word
\end{itemize}

\begin{tcolorbox}[breakable,colback=teal!4,colframe=teal!35!black,title={Before you move on}]
What does \pa{ਠਹਿਰਨਾ} mean? (\textquotedblleft{}To halt; to stop for a while\textquotedblright{}.) Say it once more.
\end{tcolorbox}
\section[\texorpdfstring{\pa{ਚੱਲਣਾ} (challṇā)}{ਚੱਲਣਾ (challṇā)}]{\pa{ਚੱਲਣਾ} (challṇā) — to set off; to go, of a bus or train}

\label{lesson:PA-C261-challna}



\begin{quote}
Before the new one: say the Punjabi for a railway platform, then the Punjabi for to halt.
\end{quote}

\subsection*{You'll want to know: \pa{ਚੱਲਣਾ}}
\textbf{\pa{ਚੱਲਣਾ}} — \emph{challṇā} — \textquotedblleft{}to set off; to go, of a bus or train\textquotedblright{}.

\begin{quote}
\pa{ਪਹਿਲੀ} \pa{ਬੱਸ} \pa{ਸਵੇਰੇ} \pa{ਛੇ} \pa{ਵਜੇ} \pa{ਚੱਲਦੀ} \pa{ਹੈ।}
\end{quote}

\begin{quote}
\emph{pahilī bass savere chhe vaje challdī hai.}
\end{quote}

\begin{quote}
The first bus sets off at six in the morning.
\end{quote}

Past (it …): \pa{ਉਹ} \pa{ਬੱਸ} \pa{ਅੱਠ} \pa{ਵਜੇ} \pa{ਚੱਲੀ}, \emph{oh bass aṭṭh vaje challī}.

\begin{cousinweb}
\textbf{\pa{ਚੱਲਣਾ}} (\emph{challṇā}) is inherited from a reconstructed \emph{*calyati}, \textquotedblleft{}goes\textquotedblright{}; CDIAL lists Punjabi \emph{callṇā}, \textquotedblleft{}to go, move\textquotedblright{}, under it (CDIAL 4721).
\end{cousinweb}

\begin{grammarlens}[title={What you've built}]
One more word for reading notices and messages.
\end{grammarlens}

\subsection*{Your turn}
\begin{itemize}
\raggedright
  \item \emph{Say it:} \emph{challṇā}
  \item \emph{Say it:} \emph{challṇā}, once more
  \item \emph{Say it:} \emph{pahilī bass savere chhe vaje challdī hai.} — The first bus sets off at six in the morning.
  \item \emph{Recall:} say the Punjabi for a railway platform, then the Punjabi for to halt, then say \emph{challṇā} again
  \item \emph{Write it:} copy \textbf{\pa{ਚੱਲਣਾ}} once, checking the whole word
\end{itemize}

\begin{tcolorbox}[breakable,colback=teal!4,colframe=teal!35!black,title={Before you move on}]
What does \pa{ਚੱਲਣਾ} mean? (\textquotedblleft{}To set off; to go, of a bus or train\textquotedblright{}.) Say it once more.
\end{tcolorbox}
\section[\texorpdfstring{\pa{ਰਵਾਨਾ} (ravānā)}{ਰਵਾਨਾ (ravānā)}]{\pa{ਰਵਾਨਾ} (ravānā) — departing; leaving, with \pa{ਹੋਣਾ}}

\label{lesson:PA-C261-ravana}



\begin{quote}
Before the new one: say the Punjabi for to halt, then the Punjabi for to set off.
\end{quote}

\subsection*{You'll want to know: \pa{ਰਵਾਨਾ}}
\textbf{\pa{ਰਵਾਨਾ}} — \emph{ravānā} — \textquotedblleft{}departing; leaving, with \pa{ਹੋਣਾ}\textquotedblright{}.

\begin{quote}
\pa{ਰੇਲ} \pa{ਸਾਢੇ} \pa{ਦਸ} \pa{ਵਜੇ} \pa{ਰਵਾਨਾ} \pa{ਹੁੰਦੀ} \pa{ਹੈ।}
\end{quote}

\begin{quote}
\emph{rel sāḍhe das vaje ravānā hundī hai.}
\end{quote}

\begin{quote}
The train leaves at half past ten.
\end{quote}

\begin{cousinweb}
\textbf{\pa{ਰਵਾਨਾ}} (\emph{ravānā}) is a loan, usually traced to Persian \emph{ravāna}, \textquotedblleft{}going, setting out\textquotedblright{}.
\end{cousinweb}

\begin{grammarlens}[title={What you've built}]
One more word for reading notices and messages.
\end{grammarlens}

\subsection*{Your turn}
\begin{itemize}
\raggedright
  \item \emph{Say it:} \emph{ravānā}
  \item \emph{Say it:} \emph{ravānā}, once more
  \item \emph{Say it:} \emph{rel sāḍhe das vaje ravānā hundī hai.} — The train leaves at half past ten.
  \item \emph{Recall:} say the Punjabi for to halt, then the Punjabi for to set off, then say \emph{ravānā} again
  \item \emph{Write it:} copy \textbf{\pa{ਰਵਾਨਾ}} once, checking the whole word
\end{itemize}

\begin{tcolorbox}[breakable,colback=teal!4,colframe=teal!35!black,title={Before you move on}]
What does \pa{ਰਵਾਨਾ} mean? (\textquotedblleft{}Departing; leaving, with \pa{ਹੋਣਾ}\textquotedblright{}.) Say it once more.
\end{tcolorbox}
